% Сборка: lualatex lab2-report.tex (рисунки — node ../segment-lab/scripts/report-figures.cjs)
\documentclass[11pt,a4paper]{article}
\usepackage[a4paper,margin=2cm,top=1.6cm,bottom=1.8cm]{geometry}
\usepackage{fontspec}
\usepackage[russian]{babel}
\usepackage{amsmath}
\usepackage{unicode-math}
\setmainfont{Liberation Serif}
\setmonofont{Liberation Mono}[Scale=MatchLowercase]
\setmathfont{TeX Gyre Termes Math}
\usepackage{titlesec}
\usepackage{enumitem}
\usepackage{microtype}
\usepackage{graphicx}
\usepackage{booktabs}
\usepackage{caption}
\usepackage{float}

\titleformat{\section}{\large\bfseries}{\thesection.}{0.5em}{}
\titleformat{\subsection}{\normalsize\bfseries}{\thesubsection.}{0.5em}{}
\titlespacing*{\section}{0pt}{1.0ex plus .3ex}{0.5ex}
\titlespacing*{\subsection}{0pt}{0.8ex plus .3ex}{0.4ex}
\setlength{\parindent}{0pt}
\setlength{\parskip}{0.45em}
\pagestyle{plain}
\captionsetup{font=small,labelsep=period,skip=4pt}
\graphicspath{{fig/}}

\usepackage[hidelinks,pdftitle={Лабораторная работа № 2. Сегментация изображений},pdfauthor={Aliaksei Charnukha}]{hyperref}

\newcommand{\fig}[2]{\begin{figure}[H]\centering\includegraphics[width=\linewidth]{#1}\caption{#2}\end{figure}}

\begin{document}

\begin{center}
  {\LARGE\bfseries Лабораторная работа № 2}\\[0.4em]
  {\large Основы обработки цифровых изображений. Вариант 14:\\
  локальная пороговая обработка; обнаружение точек, линий и перепадов яркости}
\end{center}

\section{Цель}
Реализовать два метода локальной пороговой обработки и методы обнаружения
точек, линий и перепадов яркости; подобрать изображения, на которых видны
сильные и слабые стороны методов, и сравнить результаты.

\section{Теория}

\subsection{Локальная пороговая обработка}
Изображение $f(x, y)$ с яркостями $0\text{–}255$ переводится в бинарное:
\[
  g(x, y) = \begin{cases} 1\ (\text{фон}), & f(x, y) > T(x, y),\\ 0\ (\text{объект}), & f(x, y) \le T(x, y). \end{cases}
\]
Если порог $T$ один на всё изображение, при неравномерном освещении он не
подходит одновременно светлой и тёмной частям. Поэтому при \textbf{локальной}
обработке порог считается для каждого пикселя по его окрестности. В окне
$n \times n$ с центром в $(x, y)$ находятся среднее $m$ и среднеквадратическое
отклонение (СКО) $s$:
\[
  m = \frac{1}{n^2} \sum_{\text{окно}} f, \qquad
  s = \sqrt{\frac{1}{n^2} \sum_{\text{окно}} f^2 - m^2}.
\]
СКО показывает локальный контраст: на однородном фоне $s$ мало (только шум),
на границе объекта~— велико.

\textbf{Метод Ниблэка:}
\[ T(x, y) = m(x, y) + k \cdot s(x, y), \qquad k \approx -0.2. \]

\textbf{Метод Сауволы:}
\[ T(x, y) = m(x, y) \cdot \left(1 + k \left(\frac{s(x, y)}{R} - 1\right)\right), \qquad k \approx 0.2,\ R = 128. \]
У Сауволы при $s \to 0$ порог стремится к $m(1 - k)$, то есть на пустом фоне
он заметно ниже средней яркости, и шум не попадает в объект. У Ниблэка
в той же ситуации $T \approx m$, и примерно половина пикселей шума
оказывается темнее порога.

Чтобы время не зависело от размера окна, $m$ и $s$ считаются через
интегральные изображения сумм $f$ и $f^2$: сумма по любому окну~— четыре
обращения к таблице. За границей изображение продолжается зеркально.

\subsection{Обнаружение точек, линий и перепадов}
Используется корреляция с маской $w$:
$R(x, y) = \sum_{i,j} w(i, j)\, f(x + i, y + j)$. Сумма коэффициентов всех
масок равна нулю, поэтому на участке постоянной яркости отклик нулевой.

\textbf{Точки.} Маска лапласиана с 8 в центре и $-1$ вокруг. Пиксель,
отличающийся от ровного окружения на $\Delta$, даёт $|R| = 8\Delta$, край
объекта~— не более $5\Delta$. Точка~— где $|R| \ge T$; порог задаётся долей
от $\max |R|$.

\textbf{Линии.} Четыре маски с весом $2$ вдоль направления и $-1$ по
сторонам, например для горизонтали
$\left(\begin{smallmatrix} -1 & -1 & -1 \\ 2 & 2 & 2 \\ -1 & -1 & -1 \end{smallmatrix}\right)$.
Светлая линия толщиной в пиксель с контрастом $\Delta$ даёт отклик своей
маски $6\Delta$. Для тёмных линий берётся $-R$. Режим «все направления»~— максимум
по четырём маскам.

\textbf{Перепады яркости: градиент} (первая производная). Как в лекции,
производные заменяются разностями:
\[
  G_x = f(x + 1, y) - f(x, y), \qquad G_y = f(x, y + 1) - f(x, y),
\]
а модуль градиента $|\nabla f| = \sqrt{G_x^2 + G_y^2}$ для упрощения
вычислений заменяется на $|G_x| + |G_y|$. Граница~— где $|\nabla f| \ge T$,
$T$~— доля от максимума. На вертикальной ступеньке высотой $\Delta$
$|\nabla f| = \Delta$, на диагональной~— до $2\Delta$.

\textbf{Перепады яркости: лапласиан гауссиана (LoG)} (вторая производная).
Изображение сворачивается с маской LoG $5 \times 5$ из лекции:
\[
  \begin{pmatrix}
    0 & 0 & -1 & 0 & 0 \\
    0 & -1 & -2 & -1 & 0 \\
    -1 & -2 & 16 & -2 & -1 \\
    0 & -1 & -2 & -1 & 0 \\
    0 & 0 & -1 & 0 & 0
  \end{pmatrix}.
\]
Маска объединяет сглаживание Гауссом и лапласиан. На перепаде вторая
производная меняет знак, поэтому граница~— это переход через ноль: у пары
противоположных соседей разные знаки и разность больше порога $T$.

Производные усиливают шум, поэтому во всех методах перед обработкой можно
включить сглаживание Гауссом с параметром $\sigma$.

\section{Ход работы}
Приложение написано на HTML, CSS и JavaScript и открывается в браузере без
установки. Все алгоритмы находятся в отдельном модуле \texttt{imaging.js},
интерфейс только вызывает его. Возможности:
\begin{itemize}[label=--, leftmargin=1.5em, itemsep=0.1em, topsep=0.2em]
  \item загрузка своего изображения (цветное переводится в полутона
        $Y = 0.299R + 0.587G + 0.114B$) или выбор из встроенной базы;
  \item выбор метода и параметров ползунками, предварительное сглаживание;
  \item вывод исходного изображения, карты порога $T(x, y)$ или отклика маски
        и результата, наложение результата на исходное;
  \item значения $I$, $T$ или $R$ под курсором;
  \item сравнение Ниблэка и Сауволы (или градиента и LoG) на одном
        изображении; для изображений с эталоном~— F1.
\end{itemize}

Модуль проверен 50 юнит-тестами: пороги Ниблэка и Сауволы совпадают со
scikit-image, свёртки с масками (в том числе LoG $5 \times 5$) и гауссово
сглаживание~— с OpenCV (расхождение меньше $10^{-6}$), остальное проверено на
простых примерах.

\textbf{База изображений} синтетическая, чтобы был известен правильный ответ:
4 изображения для бинаризации с эталонными масками, 2 с точками, 2 с линиями,
4 для перепадов. Качество бинаризации оценивалось F1-мерой по эталону
($\text{F1} = 2PR / (P + R)$, $P$~— точность, $R$~— полнота, объект~— чёрные пиксели).

\section{Эксперименты}

\subsection{Локальная пороговая обработка}

\begin{table}[H]
\centering\small
\caption{F1 по эталону, \%. Окно 25, Ниблэк $k = -0.2$, Саувола $k = 0.2$.
В скобках~— лучший результат при подборе окна 15–61 и $k$.}
\begin{tabular}{lcc}
\toprule
Изображение & Ниблэк & Саувола \\
\midrule
Документ с тенью          & 40.0 \ (85.5) & \textbf{100.0} \ (100.0) \\
Шум и пустые поля         & 12.8 \ (28.9) & \textbf{99.8} \ (100.0) \\
Выцветший документ        & \textbf{42.0} \ (74.9) & 0.0 \ (55.8) \\
Пятна на неровном фоне    & 41.2 \ (85.7) & \textbf{89.4} \ (100.0) \\
\bottomrule
\end{tabular}
\end{table}

\fig{doc-shadow}{Документ с тенью: исходное, Ниблэк, Саувола.}
Яркость фона меняется от светлого левого верхнего угла к тёмному правому
нижнему, поэтому один порог здесь не подходит. Оба локальных метода выделяют
текст, но Ниблэк добавляет шум на фоне между строками.

\fig{doc-noisy-margins}{Ровный фон с шумом и пустыми полями: исходное, Ниблэк, Саувола.}
Ниблэк на пустых полях даёт $T = m - 0.2\,s$, и около 40\,\% пикселей шума
становятся «текстом» (F1 = 12.8\,\%). У Сауволы на фоне $T \approx 0.8\,m$,
и фон чистый.

\fig{doc-faded}{Выцветший документ: исходное, Ниблэк (окно 61, $k = -0.8$), Саувола ($k = 0.2$).}
Это пример, где \textbf{Ниблэк лучше Сауволы}. Фон около 200, буквы темнее на
22. Саувола отнесёт букву к объекту, только если её контраст больше
$m\,k\,(1 - s/R) \approx 0.19\,m \approx 37$, поэтому весь текст пропадает.
Ниблэк отсчитывает порог от среднего по $s$ и находит бледный текст,
хотя пустая часть справа покрывается шумом.

\fig{blobs-uneven}{Пятна на неравномерном фоне: исходное, Саувола с окном 15 и 61.}
Окно должно быть заметно больше объекта. При окне 15 оно целиком попадает
внутрь пятна, порог считается по самому пятну, и пятна получаются
«дырявыми». При окне 61 F1 = 100\,\%.

\subsection{Точки}
\fig{points}{Точки с шумом: исходное, $|R|$, порог 60\,\%, порог 40\,\%.}
На изображении 10 светлых ($+100$, $|R| = 800$) и 6 тёмных ($-70$,
$|R| = 560$) точек. Плавный фон маска не замечает, края фигур дают отклик
до $5\Delta \approx 350$.

\begin{table}[H]
\centering\small
\caption{Найдено пикселей / из них настоящих точек (всего 16).}
\begin{tabular}{lcccccc}
\toprule
Порог, \% от $\max|R|$ & 30 & 40 & 50 & 60 & 70 & 90 \\
\midrule
Без шума     & 295 / 16 & 78 / 16 & 18 / 16 & \textbf{16 / 16} & 10 / 10 & 10 / 10 \\
Шум $\sigma = 6$ & 285 / 16 & 64 / 16 & 21 / 16 & \textbf{15 / 15} & 11 / 11 & 6 / 6 \\
\bottomrule
\end{tabular}
\end{table}
При пороге 60\,\% находятся все точки без ложных. При 70--90\,\% остаются
только светлые, при 30--40\,\%~— добавляются углы, края фигур и шум.

\subsection{Линии}
\fig{lines}{Линии: исходное, отклик (максимум по 4 маскам), все направления при пороге 50\,\% и 55\,\%, только горизонтальная маска (55\,\%).}
Каждая линия толщиной в пиксель выделяется своей маской ($6\Delta = 720$).
Внутри толстой полосы справа внизу отклик нулевой, а на её краях
$3\Delta = 360$~— ровно половина. Поэтому при пороге 50\,\% от полосы
остаётся контур, а при 55\,\% она уже не выделяется. Бледная линия сверху
(контраст 50 вместо 120, отклик 300) не проходит ни один из порогов~— это
недостаток порога от максимума: слабые линии теряются рядом с сильными.

\fig{wires}{Тёмные линии: исходное, горизонтальная маска в режиме «светлые», в режиме «тёмные».}
Если искать светлые линии на изображении с тёмными, маска срабатывает по
двум сторонам линии, и линия удваивается. Наклонные линии (около $75^\circ$)
не совпадают ни с одной из четырёх масок и горизонтальной маской
не выделяются.

\subsection{Перепады яркости}
Качество оценивалось по контуру, найденному на чистом изображении
(пиксель засчитывается с допуском $\pm 2$ пикселя).

\begin{table}[H]
\centering\small
\caption{F1, \%, относительно контура чистых фигур. Порог градиента 25\,\%, LoG~— 10\,\% или 20\,\% от максимума.}
\begin{tabular}{lccc}
\toprule
Метод & Чистые & Шум $\sigma = 18$ & Размытие $\sigma = 3$ \\
\midrule
Градиент                          & 91.5 & 72.7 & 84.9 \\
Градиент + Гаусс $\sigma = 1.5$   & 98.3 & \textbf{97.9} & 78.3 \\
LoG 10\,\%                        & \textbf{100.0} & 19.5 & 70.9 \\
LoG 10\,\% + Гаусс $\sigma = 1.5$ & \textbf{100.0} & 22.1 & \textbf{99.9} \\
LoG 20\,\% + Гаусс $\sigma = 3$   & 95.4 & 94.7 & 89.4 \\
\bottomrule
\end{tabular}
\end{table}

\fig{edges-clean}{Чистые фигуры: исходное, градиент (25\,\%), LoG (10\,\%).}
Порог градиента считается от максимума, а максимум даёт контрастный
треугольник (перепад 170). Граница эллипса (перепад 50) проходит порог только
на наклонных участках, где $|G_x| + |G_y|$ вдвое больше. У LoG граница~— переход
через ноль, и контуры всех фигур находятся полностью.

\fig{edges-noisy}{Фигуры с шумом: исходное, градиент, градиент после сглаживания $\sigma = 1.5$, LoG, LoG после сглаживания $\sigma = 3$ (20\,\%).}
Разности соседних пикселей реагируют на шум, и градиент даёт много ложных
точек. Лапласиан~— вторая производная, он к шуму ещё чувствительнее:
переходы через ноль покрывают весь фон. Градиенту достаточно сглаживания
$\sigma = 1.5$, LoG~— только при $\sigma = 3$ и пороге 20\,\%.

\fig{edges-blurred}{Размытые фигуры: исходное, $|\nabla f|$, градиент, LoG после сглаживания $\sigma = 1.5$.}
На размытом изображении перепад растянут на несколько пикселей, и градиент
даёт границы толщиной 3--4 пикселя (точность около 76\,\%). Переход через
ноль у LoG один, посередине перепада, поэтому контур тонкий и замкнутый.
Без сглаживания LoG находит ещё и ложные переходы через ноль на мелких
колебаниях яркости (F1 = 70.9\,\%).

\section{Результаты}
\begin{itemize}[label=--, leftmargin=1.5em, itemsep=0.2em, topsep=0.2em]
  \item Разработано веб-приложение с графическим интерфейсом. В нём реализованы
        локальная пороговая обработка методами Ниблэка и Сауволы, обнаружение
        точек и линий масками и обнаружение перепадов яркости по градиенту и
        лапласиану гауссиана.
  \item Реализация проверена 50 юнит-тестами: результаты совпадают с OpenCV и
        scikit-image с точностью до $10^{-6}$.
  \item Собрана база из 12 изображений; для четырёх изображений бинаризации
        есть эталонные маски.
  \item Бинаризация (F1): документ с тенью~— Ниблэк 40.0\,\%, Саувола
        100\,\%; пустые поля с шумом~— Ниблэк 12.8\,\%, Саувола 99.8\,\%;
        выцветший документ~— Ниблэк 42.0\,\% (74.9\,\% после подбора
        параметров), Саувола 0\,\% (55.8\,\%).
  \item Точки: при пороге 60\,\% найдены все 16 точек без ложных
        срабатываний, на зашумлённом изображении~— 15 из 16.
  \item Линии: при пороге 55\,\% найдены все линии с контрастом 120; толстая
        полоса и бледная линия с контрастом 50 не выделяются.
  \item Перепады (F1): на чистых фигурах градиент 91.5\,\%, LoG 100\,\%; на
        шуме градиент 72.7\,\% (97.9\,\% после сглаживания), LoG 19.5\,\%
        (94.7\,\% после сглаживания $\sigma = 3$); на размытых фигурах
        градиент 84.9\,\%, LoG после сглаживания 99.9\,\%.
\end{itemize}

\section{Выводы}
\begin{itemize}[label=--, leftmargin=1.5em, itemsep=0.2em, topsep=0.2em]
  \item Локальная пороговая обработка нужна при неравномерном освещении:
        один порог не может одновременно подойти светлой и тёмной частям
        изображения.
  \item Саувола лучше Ниблэка, когда на изображении есть области без
        объектов: Ниблэк превращает шум на фоне в объект. Ниблэк лучше на
        бледном малоконтрастном тексте: Саувола требует контраста примерно
        $0.2\,m$ и теряет такой текст полностью.
  \item Размер окна должен быть больше объекта, иначе объект
        бинаризуется «сам относительно себя».
  \item Маски точек и линий находят только объекты толщиной в пиксель и
        своего направления; порог от максимального отклика удобен, но
        слабые объекты теряются, если на изображении есть сильные.
  \item Градиент прост и даёт границы на чистых изображениях, но порог от
        максимума теряет слабые перепады, а на размытых границах линии
        получаются толстыми. LoG находит тонкие замкнутые контуры независимо
        от высоты перепада, но как вторая производная сильнее реагирует на
        шум и без сглаживания на зашумлённых изображениях непригоден.
\end{itemize}

\end{document}
